\documentclass[12pt,a4paper]{article}
\usepackage{fontspec}\usepackage{polyglossia}\setdefaultlanguage{russian}
\setmainfont{DejaVu Serif}\usepackage{geometry}\geometry{margin=2.3cm}
\usepackage{titlesec}\titleformat{\subsection}{\bfseries\large}{}{0pt}{}
\setlength{\parskip}{.6em}\setlength{\parindent}{0pt}
\begin{document}
{\LARGE\bfseries Речь для защиты (8--10 минут)}\\[2pt]
Исследование обобщённой дискретной модели, вариант 1, $k=4$. Ориентир: около 1000 слов с формулами, при спокойном темпе 8--9 минут; если успеваете раньше, не торопитесь, паузы на слайдах 5--7 нормальны.

\subsection*{Слайд 1. Титульный (0:20)}
Здравствуйте. Меня зовут Ирина Власенко, группа ПМИ-2.2, мой номер в списке --- четыре, поэтому $k=4$. Я исследовала вариант 1: логистическое отображение с дополнительным нелинейным слагаемым. Покажу, что происходит при росте параметра $r$ и как это сравнивается с литературой.

\subsection*{Слайд 2. Постановка задачи (1:00)}
Модель --- это $x_{n+1}=r x_n\bigl(1-\frac{x_n+4x_n^2}{5}\bigr)$. Здесь $x_n$ --- нормированная плотность популяции в поколении $n$, а $r$ --- скорость роста при малой плотности: производная в нуле равна $r$. Слагаемое в скобках описывает саморегуляцию плотностью, а $k$ задаёт, насколько нелинейно она растёт. При $k=0$ мы получаем классическое логистическое отображение. Нормировка такая, что $f(1)=0$, а максимум функции равен примерно $0{,}357\,r$. Поэтому отрезок от нуля до единицы остаётся инвариантным, пока $r$ не больше $2{,}802$. Этим и объясняется диапазон диаграммы, который мы увидим позже. Важно, что при $r$ меньше единицы популяция вымирает, а при больших $r$ начинаются колебания плотности, и именно этот переход нас интересует. Цели работы --- найти неподвижные точки и условия устойчивости, критическое $r$, цикл периода два, сценарий перехода к хаосу и сравнить всё с литературой.

\subsection*{Слайд 3. Литература (1:10)}
Я использовала три реальные работы. Первая --- статья Роберта Мэя в Nature, 1976 год, о том, что простые нелинейные рекуррентные модели дают очень сложную динамику. Вторая --- работа Мэя и Остера в American Naturalist, того же года, про бифуркации и динамическую сложность в экологических моделях. Третья --- статья Радвана 2013 года в Journal of Advanced Research об обобщённых логистических отображениях, которые сводятся к классическому. Сравнение такое. Для нашей модели из условия производная равна минус единице получается уравнение $4kx^2+3x-2(1+k)=0$. При $k=0$ оно даёт $x^*=2/3$ и $r=3$ --- классическое значение для логистического отображения. Для $k=4$ первая бифуркация наступает раньше, при $r\approx2{,}151$, но сценарий тот же. Кроме того, отношения интервалов между бифуркациями у меня получились $4{,}57$ и $4{,}64$, то есть приближаются к константе Фейгенбаума, около $4{,}669$. Конкретные значения $r$ из статей на нашу модель я напрямую не переношу.

\subsection*{Слайд 4. Неподвижные точки и устойчивость (1:20)}
Неподвижные точки: ноль и положительная точка $x^*_+=\frac{-1+\sqrt{81-80/r}}{8}$, она существует при $r>1$; второй корень отрицателен и физического смысла не имеет. Производная $f'(x)=\frac r5(5-2x-12x^2)$. В нуле она равна $r$, значит ноль устойчив при $r$ от нуля до единицы. Для положительной точки, используя уравнение неподвижной точки, производная упрощается: $1-\frac{r x^*(1+8x^*)}{5}$. Она всегда меньше единицы и убывает с ростом $r$, поэтому устойчивость теряется только через значение минус один. Приравниваем и получаем $16x^2+3x-10=0$, откуда $x_{\rm crit}=\frac{\sqrt{649}-3}{32}\approx0{,}702$, а критическое $r$ равно примерно $2{,}1511$. Итог: положительная точка устойчива при $r$ от единицы до $r_{\rm crit}$, а при $r_{\rm crit}$ происходит удвоение периода.

\subsection*{Слайд 5. Цикл периода 2 (1:20)}
Цикл --- это пара точек, которые переходят друг в друга. Они удовлетворяют уравнению $f(f(x))=x$. Здесь важный момент: функция $f$ кубическая, поэтому $f\circ f$ имеет степень девять, и после деления на $f(x)-x$ остаётся уравнение шестой степени, а не квадратное, как предполагает задание. Я проверила это в sympy. Это ещё одно подтверждение того, что к утверждениям из методички полезно относиться критически и проверять их для конкретной формулы. Но уравнение решается: вводим сумму $s$ и произведение $p$ точек цикла. Тогда $p$ выражается через $s$, а сама $s$ --- корень кубического уравнения $16rs^3+8rs^2-(19r+40)s-5(r+1)=0$. Потом $x_1$ и $x_2$ --- корни квадратного трёхчлена. Физичный цикл появляется при $r>r_{\rm crit}$. Устойчивость определяет мультипликатор $\mu=f'(x_1)f'(x_2)$; из условия $\mu=-1$ бисекцией получаем $r_2\approx2{,}4136$. Значит, цикл устойчив между $2{,}151$ и $2{,}414$.

\subsection*{Слайд 6. Численно: сходимость и цикл (1:00)}
Теперь численные эксперименты. Слева $r=0{,}9\,r_{\rm crit}$: траектория быстро сходится к неподвижной точке $0{,}6624$. Обратите внимание, что вблизи бифуркации переходный процесс идёт медленнее: амплитуда колебаний нарастает несколько десятков шагов, прежде чем устанавливается цикл. Справа $r=1{,}1\,r_{\rm crit}$: устанавливается цикл периода два с точками $0{,}532$ и $0{,}840$. В таблице --- проверка. Невязка уравнений цикла порядка $10^{-15}$. Мультипликатор при $1{,}05\,r_{\rm crit}$ равен $0{,}21$, при $1{,}1\,r_{\rm crit}$ равен $-0{,}62$: оба по модулю меньше единицы, цикл устойчив. При $1{,}2\,r_{\rm crit}$ мультипликатор $-2{,}4$ --- цикл существует, но уже неустойчив, потому что мы ушли за $r_2$.

\subsection*{Слайд 7. Бифуркации (1:00)}
Слева --- бифуркационная диаграмма. Отброшено 1500 итераций переходного процесса. Видна первая бифуркация при $2{,}151$, вторая при $2{,}414$, затем каскад удвоений и область хаоса с окном периода три около $r=2{,}66$. Каждая вертикальная линия на диаграмме --- это набор значений, которые система посещает после переходного процесса: одна точка означает равновесие, две --- цикл периода два, густая область --- сложную динамику. Диаграмма заканчивается на $2{,}8$, потому что дальше траектории уходят из отрезка $[0,1]$. Справа --- ветви цикла $x_1(r)$ и $x_2(r)$ и неподвижная точка: пунктиром показаны неустойчивые части. Следующие бифуркации я нашла методом Ньютона: $r_3\approx2{,}4710$, $r_4\approx2{,}4834$.

\subsection*{Слайд 8. Хаос и гистограмма (0:50)}
Для хаотического режима взято $r=2{,}7$, 10 тысяч итераций после отбрасывания тысячи. Одного вида графика для вывода о хаосе недостаточно, поэтому я посчитала показатель Ляпунова: он равен $0{,}494$, то есть положителен. Две траектории, различающиеся на $10^{-10}$, расходятся примерно за 40 итераций. На гистограмме есть запрещённые зоны: значения не бывают ниже $0{,}168$ и выше $0{,}964$. Пики плотности стоят на границах и в точке $0{,}428$: это образы критической точки, типичные для унимодальных отображений. В середине отрезка система бывает реже. Таким образом, гистограмма показывает, что система проводит больше всего времени вблизи границ области, а середина отрезка посещается реже, и область значений действительно ограничена с обеих сторон.

\subsection*{Слайд 9. Выводы (0:40)}
Подведу итоги. Положительная неподвижная точка устойчива при $1<r<2{,}1511$. Цикл периода два рождается в этой точке и устойчив до $2{,}4136$. Для кубического отображения цикл находится через кубическое уравнение для суммы и квадратное для точек, а не через одно квадратное. Дальше идёт каскад удвоений и хаос с положительным показателем Ляпунова. Результаты согласуются с литературой: при $k=0$ получается классическое значение $r=3$, а отношения интервалов приближаются к константе Фейгенбаума.

\subsection*{Слайд 10. Использование ИИ (0:50)}
ИИ я использовала для выкладок, кода и оформления. Ключевые запросы: найти неподвижные точки, вывести критическое $r$ из условия производной равной минус единице, написать код диаграммы и гистограммы, проверить пару точек цикла. Были две ошибки. Первая: утверждение, что уравнение цикла квадратное, --- для кубической функции оно шестой степени; я исправила через деление в sympy и редукцию к кубическому уравнению для суммы. Вторая: вывод о хаосе по виду графика; я посчитала показатель Ляпунова. ИИ помог с рутиной, а формулы, диапазоны параметра и выводы потребовали ручной проверки. Спасибо за внимание, готова ответить на вопросы.

\subsection*{Возможные вопросы}
\textbf{Почему $f'(x^*)<1$ всегда?} Потому что $f'(x^*)=1-\frac{r x^*(1+8x^*)}5$ и $x^*>0$. \textbf{Почему цикл не квадратный?} Степень $f\circ f$ равна 9, деление на $f-x$ даёт степень 6. \textbf{Чем доказан хаос?} $\lambda\approx0{,}494>0$ плюс чувствительность к $x_0$. \textbf{Почему диапазон до 2,8?} При $r>2{,}802$ максимум $f$ превышает 1, траектории выходят из $[0,1]$.
\end{document}
